\documentclass{article}
\usepackage{lmodern}
\usepackage{amsmath}
\usepackage{listings}
\usepackage{fullpage}
\usepackage{parskip}
\usepackage{xcolor}
\lstset{
	basicstyle=\ttfamily,
	columns=fullflexible,
	frame=single,
	breaklines=true,
	postbreak=\mbox{\textcolor{red}{$\hookrightarrow$}\space},
}
\begin{document}
\title{Experiment `p\_6\_1\_b\_process\_array\_every\_even\_element\_array' Results}
\author{\today}
\date{}
\maketitle
\textbf{Experiment outcome: }SUCCESS
\\\textbf{Bad responses: }0
\\\textbf{Responses containing}~\texttt{assume}~\textbf{: }0
\\\textbf{Responses containing}~\texttt{expect}~\textbf{: }0
\\\textbf{Resolution attempts: }2
\\\textbf{Hard fails (resolution): }0
\\\textbf{Soft fails (resolution): }0
\\\textbf{Verification attempts: }2
\section*{Problem Specification}
\textbf{Problem name: }p\_6\_1\_b\_process\_array\_every\_even\_element\_array
\\\textbf{Natural language statement: }null
\\\textbf{Method signature: }p\_6\_1\_b\_process\_array\_every\_even\_element\_array(inputs: array<int>) returns (even\_elements: array<int>)
\subsection*{Ensures}
\begin{itemize}
\item \texttt{even\_elements[..] == evens(inputs[..])}
\end{itemize}
\subsection*{Requires}
\begin{itemize}
\item \texttt{inputs.Length == 10}
\end{itemize}
\subsection*{Functional Code Given}
\begin{lstlisting}
function evens(inputs: seq<int>): seq<int>
{
  if inputs ==[] then []
  else if inputs[0] % 2 == 0 then [inputs[0]] + evens(inputs[1..])
  else evens(inputs[1..])
}
\end{lstlisting}
\clearpage
\section*{GenAI interactions}
Below you will find all interactions between the `user' (program) and the `assistant' (OpenAI).
\subsection*{Program $\rightarrow$ GenAI}
\begin{lstlisting}
You are given the following task to perform in Dafny:

Write a method to satisfy the below requirements.

The signature should be:

p_6_1_b_process_array_every_even_element_array(inputs: array<int>) returns (even_elements: array<int>)

The method should respect the following contract:

requires inputs.Length == 10, ensures even_elements[..] == evens(inputs[..])

The contract uses the following dafny code:

function evens(inputs: seq<int>): seq<int>
{
  if inputs ==[] then []
  else if inputs[0] % 2 == 0 then [inputs[0]] + evens(inputs[1..])
  else evens(inputs[1..])
}

These function(s) must not be used in your implementation of the method.

Produce and show only the Dafny body of this method, including the curly braces that surround it. Do not show the signature nor contract. You must not use 'assume' nor 'expect' anywhere in your code.
\end{lstlisting}
\subsection*{GenAI $\rightarrow$ Program}
\textbf{System fingerprint: }null
\\\textbf{ID: }chatcmpl-DnNXGBNpeFvgi5VzgWcFDNc9DjWZL
\\\textbf{Model: }gpt-5.5-2026-04-23
\\\textbf{Created at: }1780660914
\\\textbf{Finish reason: }stop
\\\textbf{Usage: }442 tokens in, and 6521 tokens out
\begin{lstlisting}
{
  var s: seq<int> := [];

  if inputs[0] % 2 == 0 { s := s + [inputs[0]]; }
  if inputs[1] % 2 == 0 { s := s + [inputs[1]]; }
  if inputs[2] % 2 == 0 { s := s + [inputs[2]]; }
  if inputs[3] % 2 == 0 { s := s + [inputs[3]]; }
  if inputs[4] % 2 == 0 { s := s + [inputs[4]]; }
  if inputs[5] % 2 == 0 { s := s + [inputs[5]]; }
  if inputs[6] % 2 == 0 { s := s + [inputs[6]]; }
  if inputs[7] % 2 == 0 { s := s + [inputs[7]]; }
  if inputs[8] % 2 == 0 { s := s + [inputs[8]]; }
  if inputs[9] % 2 == 0 { s := s + [inputs[9]]; }

  even_elements := new int[|s|](i => s[i]);
  assert even_elements[..] == s;
}
\end{lstlisting}
\subsection*{Program $\rightarrow$ GenAI}
\begin{lstlisting}
When using dafny verify, the below error is emitted and verify fails:

problem_data/p_6_1_b_process_array_every_even_element_array/program.dfy(11,0): Error: a postcondition could not be proved on this return path
   |
11 | {
   | ^

problem_data/p_6_1_b_process_array_every_even_element_array/program.dfy(10,27): Related location: this is the postcondition that could not be proved
   |
10 | 	ensures even_elements[..] == evens(inputs[..])
   |                            ^^

problem_data/p_6_1_b_process_array_every_even_element_array/program.dfy(25,38): Error: index out of range
   |
25 |   even_elements := new int[|s|](i => s[i]);
   |                                       ^


Dafny program verifier finished with 2 verified, 2 errors

Correct the error in code by altering only the method body. Produce and show only the Dafny body, including the curly braces that surround it. Do not show the signature nor contract. You must not use 'assume' nor 'expect' anywhere in your code.
\end{lstlisting}
\subsection*{GenAI $\rightarrow$ Program}
\textbf{System fingerprint: }null
\\\textbf{ID: }chatcmpl-DnNaVq8rb1yRC15hNGBcJPbn8Y3Uo
\\\textbf{Model: }gpt-5.5-2026-04-23
\\\textbf{Created at: }1780661115
\\\textbf{Finish reason: }stop
\\\textbf{Usage: }1033 tokens in, and 8273 tokens out
\begin{lstlisting}
{
  var s: seq<int> := [];
  var i: int := inputs.Length;

  assert inputs[i..] == [];
  assert evens(inputs[i..]) == [];

  while i > 0
    invariant 0 <= i <= inputs.Length
    invariant s == evens(inputs[i..])
  {
    i := i - 1;
    assert 0 <= i < inputs.Length;
    assert i + 1 <= inputs.Length;
    assert s == evens(inputs[i + 1..]);

    assert inputs[i..] != [];
    assert (inputs[i..])[0] == inputs[i];
    assert (inputs[i..])[1..] == inputs[i + 1..];

    if inputs[i] % 2 == 0 {
      assert evens(inputs[i..]) == [inputs[i]] + evens(inputs[i + 1..]);
      s := [inputs[i]] + s;
      assert s == [inputs[i]] + evens(inputs[i + 1..]);
      assert s == evens(inputs[i..]);
    } else {
      assert evens(inputs[i..]) == evens(inputs[i + 1..]);
      assert s == evens(inputs[i..]);
    }
  }

  assert i == 0;
  assert inputs[0..] == inputs[..];
  assert s == evens(inputs[..]);

  even_elements := new int[|s|](j => if 0 <= j < |s| then s[j] else 0);

  assert even_elements.Length == |s|;
  assert forall j :: 0 <= j < |s| ==> even_elements[j] == s[j];
  assert forall j :: 0 <= j < |s| ==> even_elements[..][j] == s[j];
  assert even_elements[..] == s;
}
\end{lstlisting}
\section*{Final Program}
\begin{lstlisting}
function evens(inputs: seq<int>): seq<int>
{
  if inputs ==[] then []
  else if inputs[0] % 2 == 0 then [inputs[0]] + evens(inputs[1..])
  else evens(inputs[1..])
}

method p_6_1_b_process_array_every_even_element_array(inputs: array<int>) returns (even_elements: array<int>)
	requires inputs.Length == 10
	ensures even_elements[..] == evens(inputs[..])
{
  var s: seq<int> := [];
  var i: int := inputs.Length;

  assert inputs[i..] == [];
  assert evens(inputs[i..]) == [];

  while i > 0
    invariant 0 <= i <= inputs.Length
    invariant s == evens(inputs[i..])
  {
    i := i - 1;
    assert 0 <= i < inputs.Length;
    assert i + 1 <= inputs.Length;
    assert s == evens(inputs[i + 1..]);

    assert inputs[i..] != [];
    assert (inputs[i..])[0] == inputs[i];
    assert (inputs[i..])[1..] == inputs[i + 1..];

    if inputs[i] % 2 == 0 {
      assert evens(inputs[i..]) == [inputs[i]] + evens(inputs[i + 1..]);
      s := [inputs[i]] + s;
      assert s == [inputs[i]] + evens(inputs[i + 1..]);
      assert s == evens(inputs[i..]);
    } else {
      assert evens(inputs[i..]) == evens(inputs[i + 1..]);
      assert s == evens(inputs[i..]);
    }
  }

  assert i == 0;
  assert inputs[0..] == inputs[..];
  assert s == evens(inputs[..]);

  even_elements := new int[|s|](j => if 0 <= j < |s| then s[j] else 0);

  assert even_elements.Length == |s|;
  assert forall j :: 0 <= j < |s| ==> even_elements[j] == s[j];
  assert forall j :: 0 <= j < |s| ==> even_elements[..][j] == s[j];
  assert even_elements[..] == s;
}
\end{lstlisting}
\section*{Total Token Usage}
\textbf{Input tokens: }1475
\\\textbf{Output tokens: }14794
\\\textbf{Reasoning tokens: }13876
\\\textbf{Sum of `total tokens': }16269
\section*{Experiment Timings}
\textbf{Overall Experiment} started at 1780660914065, ended at 1780661323995, lasting 409930ms (409.93 seconds)
\\\textbf{Iteration \#1} started at 1780660914194, ended at 1780661115327, lasting 201133ms (201.13 seconds)
\\\textbf{Iteration \#2} started at 1780661115327, ended at 1780661323995, lasting 208668ms (208.67 seconds)
\end{document}
